\section{クラス設計}\label{sec:dart-class}

\begin{table}[H]
  \centering
  \begin{tabular}{|c|l|l|} \hline
    レイヤー名                   & \multicolumn{2}{l|}{rule.dart}  \\ \hline
    レイヤー概要                  & \multicolumn{2}{p{12cm}|}{利用規約のオーバーレイをテストするためのレイヤー}           \\ \hline
    責任者                  & \multicolumn{2}{l|}{金川幸太}           \\ \hline
    \multirow{7}{*}{構成要素}       & クラス名 & Rule      \\ \cline{2-3}
                               & クラス種類 & extends StatelessWidget \\ \cline{2-3}
                               & 処理概要 & 利用規約のオーバーレイを表示する。 \\ \cline{2-3}
                               & 入力 & なし \\ \cline{2-3}
                                & 出力 & Widget \\ \cline{2-3}
                                & \multicolumn{2}{l|}{他クラス・関数との関係}  \\ \cline{2-3}
                                & \multicolumn{2}{p{12cm}|}{ボタンを押すことで、rule\_screen.dartのRuleScreen().show(context, text)を呼び出す}  \\ \hline
  \end{tabular}
\end{table}

\begin{table}[H]
  \centering
  \begin{tabular}{|c|l|l|} \hline
    レイヤー名                   & \multicolumn{2}{l|}{rule\_screen.dart}  \\ \hline
    レイヤー概要                  & \multicolumn{2}{p{12cm}|}{利用規約のオーバーレイで表示される内容を生成するコンポーネント}           \\ \hline
    責任者                  & \multicolumn{2}{l|}{金川幸太}           \\ \hline
    \multirow{7}{*}{構成要素}       & クラス名 & RuleScreen      \\ \cline{2-3}
                               & クラス種類 & OverlayEntry \\ \cline{2-3}
                               & 処理概要 & オーバーレイの内容を返す \\ \cline{2-3}
                               & 入力 & Buildcontext context, String text \\ \cline{2-3}
                                & 出力 & Material \\ \cline{2-3}
                                & \multicolumn{2}{l|}{他クラス・関数との関係}  \\ \cline{2-3}
                                & \multicolumn{2}{p{12cm}|}{ボタンを押すことで、rule\_screen.dartのRuleScreen().show(context, text)が呼び出され、オーバレイ表示
                                ElevatedButtonの処理がされると、RuleScreen().hide()によりcontrolle.closeが呼び出されオーバレイを閉じる}  \\ \hline
  \end{tabular}
\end{table}
